% Complete Section 7.3, original PDF 95-96 / printed 85-86.
% block: trimmed-polynomial-forms-heading page:95 type:section
\subsection{{The trimmed polynomial spaces of differential forms}}\label{sec:trimmed-polynomial-forms}

% block: trimmed-polynomial-definition page:95 type:ordinary-paragraph
{In view of the results just established, the complete space of polynomial $k$-forms of degree $r$ can be decomposed as}
\[
\begin{aligned}
\mathcal P_r\Lambda^k&=\mathcal P_{r-1}\Lambda^k\oplus\mathcal H_r\Lambda^k\\
&=\mathcal P_{r-1}\Lambda^k\oplus\kappa\mathcal H_{r-1}\Lambda^{k+1}\oplus d\mathcal H_{r+1}\Lambda^{k-1}.
\end{aligned}
\]
{If we drop the last summand, we get a space intermediate between $\mathcal P_{r-1}\Lambda^k$ and $\mathcal P_r\Lambda^k$, the trimmed spaces of polynomial forms of degree $r$:}
\[
\mathcal P_r^-\Lambda^k:=\mathcal P_{r-1}\Lambda^k+\kappa\mathcal H_{r-1}\Lambda^{k+1}.
\]
{Note that $\mathcal P_r^-\Lambda^0=\mathcal P_r\Lambda^0$ and $\mathcal P_r^-\Lambda^n=\mathcal P_{r-1}\Lambda^n$, but for $0<k<n$, $\mathcal P_r^-\Lambda^k$ is contained strictly between $\mathcal P_{r-1}\Lambda^k$ and $\mathcal P_r\Lambda^k$. One can compute the dimension of the dropped term $d\mathcal H_{r+1}\Lambda^{k-1}$ using the exactness of the Koszul complex and induction (see [\BookNumber{11}, Theorem 3.3]). This gives a formula for the dimension of the trimmed polynomial spaces:}
\[
\dim\mathcal P_r^-\Lambda^k=\binom{r+n}{r+k}\binom{r+k-1}k.
\]

% block: trimmed-polynomial-dimension-comparison page:96 type:ordinary-paragraph
{Comparing with \eqref{eq:complete-polynomial-form-dimension}, we see that}
\[
\dim\mathcal P_r^-\Lambda^k=\frac r{r+k}\dim\mathcal P_r\Lambda^k.
\]
{We have $d\mathcal P_r^-\Lambda^k\subset d\mathcal P_r\Lambda^k\subset\mathcal P_{r-1}\Lambda^{k+1}\subset\mathcal P_r^-\Lambda^k$. So we obtain another polynomial de Rham complex in which the spaces are trimmed:}
\begin{equation}\label{eq:trimmed-polynomial-de-rham-complex}
\mathcal P_r^-\Lambda^0\xrightarrow{d}\mathcal P_r^-\Lambda^1\xrightarrow{d}\cdots\xrightarrow{d}\mathcal P_r^-\Lambda^n\to0.
\end{equation}

% block: trimmed-polynomial-exactness page:96 type:ordinary-paragraph
{Note that in this sequence, unlike in the polynomial de Rham complex of Corollary~\ref{cor:polynomial-de-rham-exactness}, the degree $r$ stays constant. This complex is also exact. To see this, first note that if $d\kappa\omega=0$ for some differential form $\omega$, then, by the homotopy formula applied to $\kappa\omega$, we have $\kappa\omega=0$. In short, $d$ is injective on the range of $\kappa$. It follows that}
\begin{equation}\label{eq:trimmed-polynomial-kernel}
\mathcal N(d|_{\mathcal P_{r-1}\Lambda^k})=\mathcal N(d|_{\mathcal P_r^-\Lambda^k}).
\end{equation}
{Next, since $\mathcal P_r\Lambda^{k-1}=\mathcal P_r^-\Lambda^{k-1}+d\mathcal H_{r+1}\Lambda^{k-2}$, we obtain}
\begin{equation}\label{eq:trimmed-polynomial-range}
d\mathcal P_r\Lambda^{k-1}=d\mathcal P_r^-\Lambda^{k-1}.
\end{equation}
{By the exactness of the full polynomial de Rham complex, the left-hand sides of \eqref{eq:trimmed-polynomial-kernel} and \eqref{eq:trimmed-polynomial-range} are equal. Therefore, the right-hand sides are equal as well:}
\[
\mathcal N(d|_{\mathcal P_r^-\Lambda^k})=d\mathcal P_r^-\Lambda^{k-1}.
\]
{This establishes the exactness of \eqref{eq:trimmed-polynomial-de-rham-complex}.}

% block: trimmed-space-origin-independence page:96 type:remark
\Needspace{6\baselineskip}
\begin{Remark}\label{rem:trimmed-origin-independence}
{The Koszul differential was defined as the contraction of a differential form at each point $x$ with the vector in $T_x\mathbb R^n$ obtained by translating the vector in $\mathbb R^n\equiv T_0\mathbb R^n$ which points from the origin to $x$. Thus the Koszul differential depends on the choice of origin in $\mathbb R^n$. However, one can check [\BookNumber{13}, p. 331] that the polynomial space $\mathcal P_r^-\Lambda^k$ is unaffected by this choice. This observation is significant, because it allows us to define the trimmed space $\mathcal P_r^-\Lambda^k$ not only on $\mathbb R^n$ but on any affine subspace. It is then easy to see that if $\omega\in\mathcal P_r^-\Lambda^k(\mathbb R^n)$, then the trace of $\omega$ on an affine subspace $V$ belongs to $\mathcal P_r^-\Lambda^k(V)$.}
\end{Remark}
